% ---------------- Résumé en français ----------------
\chapter*{Résumé}
\addcontentsline{toc}{chapter}{Résumé en français}

La tobramycine est un aminoglycoside polycationique figurant sur la liste modèle de l'OMS des médicaments essentiels, indispensable contre \textit{Pseudomonas aeruginosa} dans la mucoviscidose et les infections nosocomiales à bacilles Gram négatif, mais disponible uniquement par voie intraveineuse et par inhalation. Bien que communément classée BCS classe~IV, l'analyse structurée des données de solubilité (94--100~mg/mL dans l'eau, soit 20--200$\times$ le critère de dose sur pH~1.2--6.8) et de perméabilité ($P_{\mathrm{app}}<1\times10^{-6}$~cm/s ; biodisponibilité orale absolue $\approx$1--2\%) démontre que la tobramycine est sans ambiguïté de \textbf{classe BCS III} : l'unique obstacle à l'absorption orale est la perméabilité. Cette reclassement redirige la problématique de formulation, de l'amélioration de la solubilité vers l'amélioration de la perméabilité, pour laquelle existent des outils validés : l'appariement ionique hydrophobe (HIP ; logP de $-2.9$ à $\approx+1.6$, soit une amélioration de 1500$\times$), les systèmes auto-émulsionnants (SEDDS), les nanoparticules, et les promoteurs d'absorption gastro-intestinale (caprate de sodium C$_{10}$, SNAC).

Sur la base de cette synthèse, un modèle pharmacocinétique de population à deux compartiments (CL = 5,5~L/h, $V_1$ = 17~L, $Q$ = 2,4~L/h, $V_2$ = 16~L, $t_{1/2}$ = 2,5~h) a été couplé à un modèle d'absorption à compartiments de transit et à une réponse pharmacodynamique de type Hill. Dix variables de formulation (logP modifié, taille particulaire, composition SEDDS, concentration en promoteur, charge polymérique, enrobage chitosane et entérique, dose) ont été optimisées par un algorithme génétique classique (population 100, 200 générations) puis par NSGA-II (population 200, 300 générations) selon quatre objectifs : maximiser $\F$, maximiser $\CMIC$, maximiser $\AMIC$ et minimiser la dose. Le front de Pareto comporte 200 formulations non dominées. Le candidat retenu (TOB-161 ; gélule entérique HIP--SEDDS 550~mg, nanoparticules de 50~nm, chitosane, 50~mM de C$_{10}$, rapport tensioactif:co-tensioactif:huile de 50:30:20) atteint $\F = 34\%$ --- une amélioration de 23$\times$ par rapport à la molécule native --- avec $C_{\max} = 9{,}5$~mg/L ($\CMIC = 9{,}5$), $\AMIC = 32$, résiduelle 0,16~mg/L et $\fmic = 30\%$.

La quantification d'incertitude (Monte Carlo, $n=200$ ; validation croisée $n=26$ ; analyse de sensibilité un facteur à la fois) identifie la dose et le logP modifié comme facteurs dominants (indices 0,42 et 0,34). Un audit méthodologique a ensuite reconstruit l'application gènes---perméabilité en un modèle d'amélioration \emph{cité, borné et émergent} : chaque facteur composant porte une plage littéraire et une incertitude graduée, le gène logP est plafonné au complexe HIP mesuré, et le multiplicateur émergent plafonne à $\times$126,3. Deux évaluations complètes et standardisées --- carte d'identité, dose unique, dosages répétés réels sur 7 jours, cibles PK/PD, population virtuelle, insuffisance rénale, effet alimentaire, sensibilité par composant, comparaison IV --- démontrent que les deux formulations candidates atteignent les cibles de pic et d'exposition à l'état d'équilibre : le gagnant exploratoire réévalué (TOB-161-A, $\times$73,9 ; 550,6 mg/j) se place dans la bande 80--120 d'AUC/MIC (87,3), et l'optimum in-loop (TOBP-001, $\times$126,3 ; 1000 mg) la dépasse avec une option 467 mg dans la bande, résiduels à la ligne de surveillance de 1 mg/L et bilan de masse vérifié à 0,1\%. Le test expérimental direct --- perméabilité Caco-2/Ussing du complexe HIP en fonction du chargement --- est défini comme premier item de la feuille de route. Le procédé de fabrication (dix opérations unitaires) est chiffré à $\approx$4,30~\$ par dose (faisabilité grade B) et la voie réglementaire 505(b)(2) est estimée à environ sept ans et 4,8~millions de dollars.

Cette thèse résout la controverse de classification BCS et fournit un cadre \textit{in silico} reproductible et entièrement traçable, aboutissant à dix candidats oraux hiérarchisés, chiffrés et évalués en risque, préparant le programme expérimental nécessaire à la mise au point du premier tobramycine oral.

\bigskip
\noindent\textbf{Mots-clés :} tobramycine ; classe BCS III ; biodisponibilité orale ; appariement ionique hydrophobe ; SEDDS ; promoteurs d'absorption ; nanoparticules ; PBPK ; algorithme génétique ; NSGA-II ; mucoviscidose.
